% ============================================================
% MODULO 06 - RESULTADOS E DISCUSSOES (parafrase propria)
% ============================================================
\section{Resultados e discuss\~oes}

A triagem convergiu para seis refer\^encias-\^ancora, distribu\'idas em tr\^es eixos: (1) verde e restaura\c{c}\~ao atencional; (2) brincar entre pares como regulador; e (3) cl\'inica humanista como ambiente facilitador.

\begin{quadro}{Quadro 1 -- S\'intese das refer\^encias-\^ancora (auditoria em 5 out. 2026)}
\renewcommand{\arraystretch}{1.25}
\begin{tabular}{p{0.8cm} p{4.2cm} p{4.2cm} p{4.6cm}}
\toprule
\textbf{N\textsuperscript{o}} & \textbf{Refer\^encia} & \textbf{Desenho} & \textbf{Achado para o caso} \\
\midrule
1 & \cite{hood2024} & Revis\~ao sistem\'atica; 458 estudos triados, 7 inclu\'idos & \'Areas com mais vegeta\c{c}\~ao e brincadeiras em ``\'arvores grandes e grama'' aparecem ligadas a menos diagn\'osticos e sintomas mais leves, mesmo ajustando renda e fatores pr\'e-natais \\
2 & \cite{kuo2004} & Estudo nacional norte-americano; N = 452 pais; 49 atividades comparadas & A mesma atividade feita em espa\c{c}o verde alivia mais os sintomas do que em \'area constru\'ida ou fechada (p $<$ 0,001); em crian\c{c}as hiperativas, s\'o o verde aberto mostrou esse efeito \\
3 & \cite{fabertaylor2009} & Experimento controlado intrasujeito; N = 17; caminhada de 20 minutos & Ap\'os caminhar em parque, as crian\c{c}as pontuaram melhor em testes objetivos de aten\c{c}\~ao do que ap\'os percurso urbano equivalente; cada participante serviu de pr\'oprio controle \\
4 & \cite{taylor2011} & Survey nacional; N = 421; local habitual de brincar & Quem brinca com frequ\^encia em verde apresenta sintomas mais leves, em todos os sexos e faixas de renda; os autores prop\~oem ensaios de ``dose verde'' regular \\
5 & \cite{damasceno2025} & Interven\c{c}\~ao de 6 meses; grupo interven\c{c}\~ao \textit{versus} controle; SNAP-IV, guias e di\'arios; Nordeste do Brasil & O tipo de TDAH se manteve, mas a intensidade recuou de ``demais'' para ``bastante/pouco''; melhoraram a agita\c{c}\~ao constante, a espera de turno e a aceita\c{c}\~ao de regras, com mais calma e abertura social \\
6 & \cite{yang2019} & Transversal; N = 59.754; China; \'indice de vegeta\c{c}\~ao a 500 m da escola & Cada 0,1 a mais no \'indice de verde do entorno escolar correspondeu a chance 13\% menor de sintomas de TDAH \\
\bottomrule
\end{tabular}
\par\smallskip
{Fonte: dados da revis\~ao; elaborado pelas autoras (2026).}
\end{quadro}

\subsection{Eixo 1 -- Verde e restaura\c{c}\~ao atencional}

H\'a converg\^encia na literatura, apesar dos m\'etodos distintos: \'areas arborizadas, quintais, parques e caminhos verdes est\~ao associados a melhor concentra\c{c}\~ao logo ap\'os a exposi\c{c}\~ao e a queixas menos intensas por parte de pais e professores \cite{hood2024,yang2019}. O \'unico experimento controlado identificado com crian\c{c}as com TDAH mediu aten\c{c}\~ao antes e depois de caminhadas curtas: o desempenho foi consistentemente superior quando o trajeto atravessava um parque, e n\~ao ruas bem cuidadas \cite{fabertaylor2009}. A Teoria da Restaura\c{c}\~ao da Aten\c{c}\~ao oferece a explica\c{c}\~ao mais aceita: a natureza prende a aten\c{c}\~ao de forma espont\^anea, deixando descansar o sistema atencional que o TDAH sobrecarrega \cite{kuo2004,taylor2011}. Com honestidade: as escalas variam (SNAP-IV, SDQ, testes objetivos), h\'a poucos experimentos com pr\'e-escolares, e nada disso sugere cura -- apenas um contexto restaurador que se soma ao tratamento.

\subsection{Eixo 2 -- Brincar entre pares como regulador}

O brincar \'e a principal ocupa\c{c}\~ao da inf\^ancia e funciona como laborat\'orio de conviv\^encia: nele a crian\c{c}a experimenta pertencimento sem exclus\~ao, lida com a frustra\c{c}\~ao e exercita esperar, planejar e sustentar o interesse \cite{figueiredo2015}. Revis\~oes nacionais sobre o brincar em crian\c{c}as com TDAH destacam ganhos de expressividade, enfrentamento de problemas e aten\c{c}\~ao \cite{damascenoeduc2025}. A interven\c{c}\~ao brasileira de seis meses registrou redu\c{c}\~ao em itens como n\~ao seguir instru\c{c}\~oes, dispersar-se com est\'imulos externos, falar em excesso e resistir a brincadeiras calmas, com o aparecimento de estados de tranquilidade e maior abertura ao contato social \cite{damasceno2025}. O brincar de rua soma tr\^es ingredientes que dialogam diretamente com a hiperatividade: regras que as pr\'oprias crian\c{c}as negociam, corpo em movimento amplo e menos tempo diante de telas.

\subsection{Eixo 3 -- Cl\'inica humanista como ambiente facilitador}

A tradi\c{c}\~ao humanista compreende a ludoterapia como rela\c{c}\~ao viva, e n\~ao como t\'ecnica sobre a crian\c{c}a: o v\'inculo intersubjetivo que nasce na sess\~ao transborda para a vida cotidiana, com corpo presente, autoria e autoaceita\c{c}\~ao \cite{axline1984,rogers2020}. Quando, fora do setting, a crian\c{c}a encontra um contexto que a recebe -- pares que a chamam pelo nome, uma rua que comporta sua energia --, ela chega \`a sess\~ao menos defensiva e mais capaz de usar aquele espa\c{c}o para integrar o vivido, em vez de apenas descarreg\'a-lo. \'E isso que torna compreens\'ivel a tranquiliza\c{c}\~ao vista ao mesmo tempo em casa e na cl\'inica, sem que a terapeuta tenha alterado sua forma de trabalhar.

\begin{quadro}{Quadro 2 -- M., 5 anos, antes e depois da mudan\c{c}a de endere\c{c}o (registro cl\'inico-materno, observacional)}
\renewcommand{\arraystretch}{1.25}
\begin{tabular}{p{3.4cm} p{4.4cm} p{4.4cm} p{2.2cm}}
\toprule
\textbf{Dom\'inio} & \textbf{Antes (contexto 1)} & \textbf{Depois (contexto 2; 4--8 semanas)} & \textbf{Fonte} \\
\midrule
Aten\c{c}\~ao sustentada & N\~ao permanecia diante da televis\~ao; levantava e mexia em objetos & Mant\'em brincadeiras calmas por 10--15 minutos; desenho e jogo & M\~ae; di\'ario de campo \\
Hiperatividade & Agita\c{c}\~ao cont\'inua; ``n\~ao para'' & Corre e brinca, mas volta \`a calma; intervalos sentada & M\~ae; sess\~oes \\
Intera\c{c}\~ao social & Poucos pares; brincar solit\'ario e interno & Brinca todo dia na vizinhan\c{c}a; \'e chamada pelas colegas & M\~ae \\
Sess\~ao terap\^eutica & Troca r\'apida de brinquedos; fala acelerada & Permanece mais tempo; espera o turno; pede com palavras & Terapeuta \\
\bottomrule
\end{tabular}
\par\smallskip
{Fonte: prontu\'ario e relato materno (2026). Dados observacionais, sem escala padronizada.}
\end{quadro}
